\chapter*{Resumen}
\addcontentsline{toc}{chapter}{Resumen}

La función de respuesta hemodinámica (HRF) describe la transformación temporal entre una perturbación neuronal y la señal BOLD observada mediante resonancia magnética funcional (fMRI). En la práctica, la forma, amplitud y latencia de la HRF varían entre sujetos, regiones cerebrales y condiciones experimentales. Esta variabilidad puede sesgar la estimación de activación cuando se utiliza una HRF canónica fija. El presente trabajo desarrolló una red neuronal informada por la física (PINN) para estimar una HRF específica de la señal y parámetros efectivos del modelo Balloon--Windkessel, combinando fidelidad a los datos con restricciones dadas por un sistema de ecuaciones diferenciales fisiológicas.

La contribución metodológica principal no consistió únicamente en entrenar una PINN. Durante el desarrollo se comprobó que una red global en el tiempo podía alcanzar residuos físicos bajos y, aun así, generar respuestas diferentes para bloques de estímulo repetidos. A partir de este diagnóstico se reformuló el método como una \textbf{PINN inversa en una ventana de entrenamiento}: la red infiere los estados latentes y los parámetros $\epsilon$, $\tau$ y $\alpha$; posteriormente, la predicción del bloque retenido se obtiene mediante una integración numérica independiente del modelo Balloon--Windkessel. Este procedimiento separa la inferencia de parámetros de la evaluación fuera de muestra y reduce la posibilidad de que la red memorice una trayectoria temporal particular.

La metodología se verificó primero en un entorno sintético con ground truth conocido. Se generaron cuatro escenarios temporales, tres niveles de relación señal--ruido (SNR 10, 5 y 2) y cinco réplicas por escenario, totalizando 60 experimentos para cada modelo. El ruido combinó componentes blanca, autorregresiva AR(1) y deriva lenta. Se compararon tres enfoques: un modelo lineal general (GLM) con HRF canónica, una red multicapa puramente basada en datos (MLP) y la PINN inversa propuesta. La PINN alcanzó un RMSE BOLD medio de $0{,}0577\%$, $0{,}0655\%$ y $0{,}0800\%$ para SNR 10, 5 y 2, respectivamente, con $R^2$ fuera de muestra superior a $0{,}98$ en los tres niveles. La HRF operacional de un evento de un segundo se recuperó con RMSE medio entre $0{,}0124\%$ y $0{,}0181\%$ BOLD. Las diferencias frente al GLM y la MLP fueron significativas en la verificación sintética. Debido a que el modelo generador coincide con el modelo incorporado en la PINN, estos resultados deben interpretarse como una verificación controlada bajo especificación correcta y no como una demostración completa de robustez frente a discrepancias fisiológicas.

La aplicación real utilizó datos de tarea motora 3T del Human Connectome Project, sujeto 100206, corridas LR y RL, y regiones M1 izquierda y derecha. Se extrajeron señales BOLD regionales, se modelaron movimientos, condiciones no objetivo y deriva temporal, y se evaluó la generalización entre los dos bloques de cada escenario en ambas direcciones. En las ocho evaluaciones direccionales, la PINN presentó el menor RMSE medio ($0{,}5019\%$ BOLD) y ganó cuatro casos, frente a dos del GLM y dos de la MLP. No se observaron diferencias globales significativas en RMSE, MAE o $R^2$. La HRF estimada por la PINN mostró alta consistencia de forma entre bloques ($r$ medio $=0{,}958$), pero la estabilidad de los parámetros fue desigual: $\alpha$ presentó una diferencia simétrica media de $11{,}17\%$, mientras $\epsilon$ y $\tau$ alcanzaron $34{,}97\%$ y $40{,}48\%$. Las pruebas inferenciales reales se consideran exploratorias, puesto que las dos direcciones de cada escenario comparten el mismo par de bloques y todos los escenarios provienen de un único sujeto.

En conjunto, el trabajo demuestra una implementación sustantiva de un problema inverso neurovascular, una verificación sintética extensa, una aplicación real reproducible y una discusión explícita de identificabilidad y alcance. Los resultados respaldan una prueba de concepto metodológica: la estructura fisiológica puede actuar como un regularizador útil y entregar HRF interpretables bajo condiciones controladas, mientras que la validación real requiere más sujetos, una ablación física directa y experimentos con desajuste de modelo antes de formular conclusiones poblacionales o clínicas.

\textbf{Palabras clave:} fMRI, BOLD, función de respuesta hemodinámica, PINN, Balloon--Windkessel, problema inverso, identificabilidad, Human Connectome Project.

\chapter*{Abstract}
\addcontentsline{toc}{chapter}{Abstract}

The hemodynamic response function (HRF) describes the temporal transformation between neural perturbation and the BOLD signal measured by functional magnetic resonance imaging. Its shape, amplitude, and latency vary across subjects, regions, and experimental conditions, potentially biasing activation estimates obtained with a fixed canonical HRF. This work developed a physics-informed neural network (PINN) to estimate signal-specific HRFs and effective Balloon--Windkessel parameters by combining BOLD data fidelity with physiological ordinary differential equation constraints.

The main methodological contribution was a reformulation from a global soft-physics trajectory network to an inverse-window PINN followed by independent ODE integration. The network infers latent states and the parameters $\epsilon$, $\tau$, and $\alpha$ from a training block, whereas the held-out response is predicted by numerically integrating the physiological model. This design provides a stricter out-of-sample test than directly evaluating the neural network trajectory.

Controlled verification included four temporal scenarios, three noise levels, and five replicates, totaling 60 synthetic experiments per model. The proposed approach was compared against a canonical-HRF GLM and a data-only MLP. Across SNR 10, 5, and 2, the PINN achieved mean held-out BOLD RMSE values of $0.0577\%$, $0.0655\%$, and $0.0800\%$, respectively, with held-out $R^2$ above $0.98$. Operational one-second-event HRFs were recovered with mean RMSE between $0.0124\%$ and $0.0181\%$ BOLD. Since synthetic data were generated with the same physiological model used for inference, these results constitute controlled verification under correct model specification rather than a complete robustness validation.

Real-data evaluation used motor-task fMRI from Human Connectome Project subject 100206, two phase-encoding runs, and bilateral primary motor cortex signals. The PINN obtained the lowest mean held-out RMSE ($0.5019\%$ BOLD) and won four of eight directional comparisons, but RMSE, MAE, and $R^2$ differences were not statistically conclusive. Estimated HRF shapes were highly consistent across blocks (mean $r=0.958$), whereas parameter identifiability was unequal, particularly for $\tau$. Real-data inference was treated as exploratory because bidirectional evaluations are dependent and all observations came from one subject.

The resulting project is a reproducible methodological proof of concept, combining physiological modeling, inverse learning, extensive synthetic experiments, real fMRI application, statistical analysis, and explicit limitations. Population-level or clinical validation will require larger cohorts, matched physics ablations, and model-mismatch experiments.
